\section{もとの積分が戻ったら等式を解く}\label{節:部分積分の循環}

指数関数と三角関数の積では，整式のように微分を繰り返して $0$ になる因子がない。その代わり，正弦と余弦は微分によって互いに移り変わる。二回の部分積分で，もとの積分が再び現れることがある。そこで計算を失敗と決めず，もとの積分を未知数として含む等式が得られていないかを調べる。

\begin{example}{指数関数と正弦の積}
$\int\e^x\sin x\dd x$ を求めよ。
\end{example}

積分定数に惑わされずに等式を整理するため，まず
\[
 I(x)=\int_0^x\e^t\sin t\dd t,\qquad
 J(x)=\int_0^x\e^t\cos t\dd t
\]
と定める。$I,J$ は $x$ の関数であり，基本定理から $I'(x)=\e^x\sin x$ である。したがって $I(x)$ が求まれば，それに任意定数を加えて不定積分の答えにできる。

第一の積分で $\sin t$ を微分し，$\e^t$ を積分すると，
\[
\begin{aligned}
 I(x)
 &=[\e^t\sin t]_0^x-\int_0^x\e^t\cos t\dd t\\
 &=\e^x\sin x-J(x).
\end{aligned}
\]
続いて $J(x)$ でも $\e^t$ を積分する側に選ぶ。$\cos t$ の微分は $-\sin t$ だから，
\[
\begin{aligned}
 J(x)
 &=[\e^t\cos t]_0^x-\int_0^x\e^t(-\sin t)\dd t\\
 &=\e^x\cos x-1+I(x).
\end{aligned}
\]
ここで $I(x)$ が戻ってきた。第二式を第一式へ代入すると
\[
\begin{aligned}
 I(x)
 &=\e^x\sin x-\{\e^x\cos x-1+I(x)\},\\
 2I(x)&=\e^x(\sin x-\cos x)+1.
\end{aligned}
\]
よって $I(x)=\frac12\e^x(\sin x-\cos x)+\frac12$ である。最後の $\frac12$ も任意定数へまとめれば，
\[
 \boxed{\int\e^x\sin x\dd x
 =\frac12\e^x(\sin x-\cos x)+C}
\]
を得る。微分による確認は
\[
 \frac12\e^x(\sin x-\cos x)
 +\frac12\e^x(\cos x+\sin x)
 =\e^x\sin x
\]
である。

ここで役立ったのは，「同じ積分が戻る」という事実だけではない。右辺へ $-I(x)$ として戻ったので，移項すると $2I(x)$ が得られたことが重要である。もし同じ式を単に逆向きにたどり，$I(x)=I(x)$ に戻っただけなら，新しい情報はない。循環した後には，何倍の $I(x)$ が残ったのかを確かめる。

不定積分のまま同じ計算を進めることもできる。その場合は，途中の等式が任意定数の違いを許していることを忘れず，最後に一つの $C$ にまとめる。定積分として $I,J$ を定めたのは，等式の両辺を同じ関数として固定し，移項の意味をはっきりさせるためである。

